\documentclass[11pt,a4paper]{article}

\usepackage[utf8]{inputenc}
\usepackage[T1]{fontenc}
\usepackage[spanish]{babel}
\usepackage{amsmath,amssymb}
\usepackage[margin=2.5cm]{geometry}
\usepackage{hyperref}

\newcommand{\code}[1]{\texttt{#1}}
\emergencystretch=2em

\title{SpatFD v2: diseño de la predicción y funcionamiento del \emph{solver}}
\author{Documento de diseño}
\date{Octubre de 2026}

\begin{document}
\maketitle

\noindent
Este documento describe lo que se diseñó para la predicción en SpatFD v2. El
centro es el \emph{solver} de kriging: qué sistema resuelve, cómo lo resuelve
y cómo comprobar a mano que sus soluciones son correctas. La última parte
plantea cómo se integraría un método \code{lambda} multivariado, que todavía
\textbf{no está implementado}.

\tableofcontents

% ======================================================================
\section{Panorama del diseño}

\subsection{El objeto y el registro de modelos}

Un objeto \code{SpatFD} guarda los \emph{scores} de los componentes
principales de cada variable y, como atributos, lo necesario para volver de
los \emph{scores} a las curvas. Los modelos de variograma se registran dentro
del mismo objeto:

\begin{itemize}
  \item \code{add\_model()} traduce el modelo (si viene de \code{gstat}), lo
        compila y lo guarda con un nombre. La compilación convierte cada
        variograma en una función de la distancia, $\gamma(h)$, y se hace una
        sola vez.
  \item Cada modelo registrado guarda tres cosas: la tabla original, los
        variogramas compilados y la meseta total de cada uno (la suma de sus
        mesetas parciales).
  \item \code{sfd\$models} es la tabla resumen de lo registrado. Se calcula al
        consultarla.
\end{itemize}

\subsection{Las capas de la predicción}

\code{predict()} es deliberadamente delgada. El trabajo está repartido así:

\begin{center}
\begin{tabular}{ll}
\hline
Función & Qué hace \\
\hline
\code{prediction\_setup()} & Verifica argumentos y prepara sitios, valores y modelo \\
\code{score\_covariances()} & Convierte los variogramas en funciones de covarianza \\
\code{simple\_kriging()} & El \emph{solver}: predice los \emph{scores} y su varianza \\
\code{krige\_scores()} & Une las tres anteriores \\
\code{predict\_scores()} & Reconstruye las curvas a partir de los \emph{scores} \\
\code{predict\_lambda()} & Método \code{lambda}, con su propio sistema \\
\code{predict()} & Escoge el método y arma el resultado \\
\hline
\end{tabular}
\end{center}

Los dos métodos devuelven lo mismo: para cada variable, las curvas predichas
(\code{\$pred}) y su varianza de predicción (\code{\$var}). Los gráficos, los
mapas y la validación cruzada (\code{leave\_one\_out()}) solo dependen de esa
salida, no del método.

% ======================================================================
\section{De variogramas a covarianzas}

El kriging simple trabaja con covarianzas. De cada variograma compilado se
obtiene
\[
  C(h) = C(0) - \gamma(h), \qquad C(0) = \text{suma de las mesetas parciales}.
\]
Tres detalles importan para verificar resultados:
\begin{itemize}
  \item El efecto pepita vale cero en $h = 0$, de modo que $C(0)$ incluye la
        pepita y $C(h)$ para $h > 0$ no. Por eso la predicción en un sitio
        observado devuelve el valor observado.
  \item Para un variograma cruzado, $C(0)$ es la suma de sus mesetas
        parciales, que pueden ser negativas. Se supone covarianza cruzada
        simétrica: $C_{ij}(h) = C_{ji}(h)$.
  \item Los modelos sin meseta no tienen covarianza y se rechazan.
\end{itemize}

% ======================================================================
\section{El \emph{solver}}

\subsection{El problema}

Un \emph{campo} es un \emph{score}, por ejemplo \code{PM10.PC1}. Hay $q$
campos; el campo $j$ se observó en $n_j$ sitios con valores $z_j$. Se apilan
todas las observaciones en un vector
\[
  z = \begin{bmatrix} z_1 \\ \vdots \\ z_q \end{bmatrix}
  \in \mathbb{R}^{N}, \qquad N = n_1 + \dots + n_q .
\]
El predictor de kriging simple con media cero para el campo $j$ en un sitio
nuevo $s_0$ y su varianza son
\begin{equation}
  \hat f_j(s_0) = c_{0j}^{T}\,\Sigma^{-1} z,
  \qquad
  \sigma_j^2(s_0) = C_{jj}(0) - c_{0j}^{T}\,\Sigma^{-1} c_{0j}.
  \label{eq:kriging}
\end{equation}

\subsection{La matriz por bloques}

$\Sigma$ es la covarianza entre todas las observaciones. Se arma por bloques,
uno por cada pareja de campos: el bloque $(i,j)$ tiene tamaño $n_i \times
n_j$ y es la covarianza $C_{ij}$ evaluada en las distancias entre los sitios
del campo $i$ y los del campo $j$. Con dos \emph{scores} de una variable $a$ y
uno de una variable $b$:
\[
  \Sigma =
  \begin{bmatrix}
    \Sigma_{a_1 a_1} & \Sigma_{a_1 a_2} & \Sigma_{a_1 b_1} \\
    \Sigma_{a_2 a_1} & \Sigma_{a_2 a_2} & \Sigma_{a_2 b_1} \\
    \Sigma_{b_1 a_1} & \Sigma_{b_1 a_2} & \Sigma_{b_1 b_1}
  \end{bmatrix},
  \qquad
  c_{0,a_1} =
  \begin{bmatrix}
    C_{a_1 a_1}(d_{a,0}) \\ C_{a_2 a_1}(d_{a,0}) \\ C_{b_1 a_1}(d_{b,0})
  \end{bmatrix},
\]
donde $d_{a,0}$ son las distancias de los sitios de $a$ al sitio nuevo.

Si el modelo no tiene variograma cruzado para una pareja, su bloque es cero.
De ahí sale la unificación de los dos casos:
\begin{itemize}
  \item \textbf{Sin variogramas cruzados}, $\Sigma$ es diagonal por bloques y
        el sistema se separa en un kriging independiente por \emph{score}.
  \item \textbf{Con variogramas cruzados}, es el sistema de cokriging simple.
\end{itemize}
Los sitios no tienen que coincidir entre variables: cada bloque usa los
sitios de sus dos campos.

\subsection{El algoritmo}

\begin{enumerate}
  \item \textbf{Distancias.} Los \emph{scores} de una misma variable comparten
        sitios, así que cada matriz de distancias se calcula una vez por
        conjunto distinto de sitios. Las distancias se calculan coordenada a
        coordenada para que la de un sitio consigo mismo sea exactamente
        cero, de lo que depende el efecto pepita.
  \item \textbf{Grupos.} Dos campos pertenecen al mismo grupo si están
        conectados por covarianzas cruzadas, directamente o a través de otros
        campos. Los grupos no comparten información y se resuelven por
        aparte; el resultado es idéntico al del sistema completo.
  \item \textbf{Armado.} Para cada grupo se construyen $\Sigma$ y la matriz
        $C_0$, cuyas columnas son los $c_{0j}$ de todos los sitios nuevos y
        todos los campos. Como $\Sigma$ es simétrica, solo se evalúa un bloque
        de cada pareja y el otro es su traspuesta.
  \item \textbf{Factorización.} Se calcula la factorización de Cholesky
        $\Sigma = R^{T} R$, con $R$ triangular superior, una sola vez por
        grupo. Si falla, $\Sigma$ no es definida positiva (sitios duplicados o
        modelo inválido) y la función se detiene con ese mensaje.
  \item \textbf{Solución.} Se resuelven dos sistemas triangulares,
        \[
          V = R^{-T} C_0, \qquad w = R^{-T} z,
        \]
        y con ellos
        \begin{equation}
          \hat f = V^{T} w,
          \qquad
          \sigma^2 = C(0) - \text{(suma por columnas de } V^2\text{)}.
          \label{eq:triangular}
        \end{equation}
\end{enumerate}

\paragraph{Por qué \eqref{eq:triangular} equivale a \eqref{eq:kriging}.}
Como $\Sigma^{-1} = R^{-1} R^{-T}$,
\[
  C_0^{T}\,\Sigma^{-1} z = (R^{-T} C_0)^{T} (R^{-T} z) = V^{T} w,
  \qquad
  c_{0}^{T}\,\Sigma^{-1} c_{0} = \lVert R^{-T} c_{0} \rVert^2 ,
\]
y esta última es la suma de los cuadrados de la columna correspondiente de
$V$. Los pesos de kriging $\Sigma^{-1} c_0$ no se calculan explícitamente:
basta una resolución triangular en lugar de dos.

\subsection{De los \emph{scores} a las curvas}

Con los \emph{scores} predichos de una variable, la curva es
\[
  \hat\chi_{s_0}(t) = \mu(t) + \sum_{k=1}^{K} \hat f_k(s_0)\,\xi_k(t),
\]
y se calcula sobre los coeficientes de la base: los coeficientes de la media
más los de las funciones propias multiplicados por los \emph{scores}. La
varianza reportada es $\sum_k \sigma_k^2(s_0)$.

% ======================================================================
\section{Cómo verificar las soluciones}

\subsection{A mano, con un solo \emph{score}}

El siguiente código resuelve el sistema directamente con \code{solve()} y lo
compara con el \emph{solver}. Usa los datos de México y un modelo fijo.

{\small
\begin{verbatim}
library(SpatFD)
data(air_mexico)

pm10 <- spatfd(pm10_mexico, coords = as.matrix(coords_pm10_mexico),
               basis = "fourier", nbasis = 21, lambda = 0.000001,
               nharm = 2, name = "PM10")
add_model(pm10, list(
  PM10.PC1 = gstat::vgm(200000, "Exp", 8000, nugget = 20000),
  PM10.PC2 = gstat::vgm(60000,  "Sph", 15000)), name = "m")
nuevo <- cbind(X = 490000, Y = 2150000)

# --- A mano ---
modelo <- attr(pm10, "models")[["m"]]
sitios <- pm10$scores[, c("X", "Y")]
z      <- pm10$scores[, "PM10.PC1"]

covarianza <- function(h)
  modelo$sills[["PM10.PC1"]] - modelo$variograms$PM10.PC1(h)

Sigma <- covarianza(as.matrix(dist(sitios)))
c0    <- covarianza(sqrt(colSums((t(sitios) - as.vector(nuevo))^2)))

pesos <- solve(Sigma, c0)
c(prediccion = sum(pesos * z),
  varianza   = covarianza(0) - sum(pesos * c0))

# --- Con el paquete ---
SpatFD:::krige_scores(pm10, nuevo, model = "m")
\end{verbatim}
}

\noindent
Los dos caminos dan lo mismo: predicción $273.0496$ y varianza $132355.0234$
para \code{PM10.PC1}.

\subsection{Propiedades que toda solución debe cumplir}

\begin{itemize}
  \item \textbf{Interpolación exacta.} Al predecir en un sitio observado se
        recupera el \emph{score} observado y la varianza es cero. En el
        ejemplo, la diferencia es del orden de $10^{-14}$ y la varianza de
        $10^{-11}$.
  \item \textbf{Varianza acotada.} Con el método \code{scores}, $0 \le
        \sigma_j^2(s_0) \le C_{jj}(0)$. Lejos de todos los sitios la varianza
        tiende a la meseta y la predicción a cero, es decir, la curva tiende a
        la media.
  \item \textbf{Bloques independientes.} Con un modelo sin variogramas
        cruzados, predecir varias variables juntas da lo mismo que predecir
        cada una por aparte.
  \item \textbf{Matriz definida positiva.} Los valores propios de $\Sigma$
        deben ser positivos. En el ejemplo el menor es $91182.6$.
  \item \textbf{Simetría.} Intercambiar el orden de las variables no cambia
        el resultado.
\end{itemize}

\subsection{Comparaciones externas}

\begin{itemize}
  \item \textbf{Contra gstat.} El script \code{dev/02\_comparacion\_gstat.R}
        compara variogramas, kriging, cokriging y \code{lambda}. Las
        diferencias relativas están entre $10^{-16}$ y $10^{-13}$.
  \item \textbf{Contra la versión CRAN.} Las curvas y varianzas de los
        métodos \code{scores} y \code{lambda} coinciden con las de
        \code{KS\_scores\_lambdas()} con diferencias del orden de $10^{-16}$.
  \item \textbf{Pruebas automáticas.} Los archivos de \code{tests/testthat}
        repiten estas comparaciones con datos simulados, incluido un caso de
        cokriging con sitios distintos por variable, pepita y mesetas
        cruzadas negativas.
\end{itemize}

Una diferencia mayor que $10^{-8}$ en cualquiera de estas comparaciones
indica un error, no redondeo.

% ======================================================================
\section{El método \code{lambda} univariado}

\subsection{El sistema}

La curva se predice como combinación lineal de las curvas observadas, con un
peso escalar por sitio:
\[
  \hat\chi_{s_0}(t) = \sum_{i=1}^{n} \lambda_i\,\chi_{s_i}(t).
\]
Los pesos minimizan $E\lVert \chi_{s_0} - \hat\chi_{s_0}\rVert^2$. Si las
curvas se escriben con sus componentes principales, el producto interno entre
dos curvas centradas es
\[
  E\langle \chi_{s_i}, \chi_{s_j} \rangle
  = \sum_{k}\sum_{l} \operatorname{Cov}\bigl(f_k(s_i), f_l(s_j)\bigr)
    \langle \xi_k, \xi_l \rangle
  = \sum_{k} C_{kk}(s_i - s_j),
\]
porque las funciones propias son ortonormales. Por eso el sistema es
\[
  \Omega\,\lambda = c,
  \qquad
  \Omega = \sum_{k=1}^{K} \Omega_k, \quad c = \sum_{k=1}^{K} c_k ,
\]
con $\Omega_k$ y $c_k$ las matrices de covarianza del \emph{score} $k$. Una
consecuencia: los variogramas cruzados entre \emph{scores} de la misma
variable no intervienen en \code{lambda}, porque $\langle \xi_k, \xi_l
\rangle = 0$ para $k \ne l$.

\subsection{Lo que se implementó}

Por decisión explícita, igual que la versión CRAN:
\begin{itemize}
  \item Los pesos se aplican a las curvas \emph{sin centrar}. En el ejemplo
        de la sección anterior los pesos suman $0.844$, no uno, así que la
        curva predicha queda encogida hacia cero.
  \item La varianza es $\sum_k \left(\eta_k - 2\,c_k^{T}\lambda +
        \lambda^{T}\Omega_k\lambda\right)$, con $\eta_k$ el valor propio del
        componente y no la meseta del modelo. Puede salir negativa.
  \item Solo una variable y modelos sin variogramas cruzados.
\end{itemize}

\subsection{Cómo verificarlo}

Los pesos y la curva se reproducen a mano, continuando el código anterior:

{\small
\begin{verbatim}
cov2  <- function(h)
  modelo$sills[["PM10.PC2"]] - modelo$variograms$PM10.PC2(h)
D     <- as.matrix(dist(sitios))
d0    <- sqrt(colSums((t(sitios) - as.vector(nuevo))^2))

lambda <- solve(covarianza(D) + cov2(D), covarianza(d0) + cov2(d0))
curva  <- attr(pm10, "fd")$PM10$coefs[, rownames(sitios)] %*% lambda

pred <- predict(pm10, nuevo, model = "m", method = "lambda")
max(abs(curva - pred$PM10$pred$coefs))      # del orden de 1e-13
\end{verbatim}
}

% ======================================================================
\section{Integración de \code{lambda} multivariado}

\noindent
\textbf{Estado: no implementado.} Lo que sigue es una propuesta. Es la
extensión directa del criterio univariado y no existe en la versión CRAN, de
modo que no hay una implementación de referencia contra la cual compararla.
Conviene contrastarla con la formulación de Bohórquez, Giraldo y Mateu (2017)
antes de programarla.

\subsection{El sistema}

Con dos variables $a$ y $b$, el predictor de la curva de $a$ usaría también
las curvas de $b$:
\[
  \hat\chi^{a}_{s_0}(t) = \mu^{a}(t)
  + \sum_{i=1}^{n_a} \lambda_i \bigl(\chi^{a}_{s_i}(t) - \mu^{a}(t)\bigr)
  + \sum_{j=1}^{n_b} \nu_j \bigl(\chi^{b}_{s_j}(t) - \mu^{b}(t)\bigr).
\]
Minimizar $E\lVert \chi^{a}_{s_0} - \hat\chi^{a}_{s_0}\rVert^2$ lleva a un
sistema por bloques, ahora con un bloque por \emph{variable}:
\begin{equation}
  \begin{bmatrix}
    \Omega_{aa} & \Omega_{ab} \\
    \Omega_{ba} & \Omega_{bb}
  \end{bmatrix}
  \begin{bmatrix} \lambda \\ \nu \end{bmatrix}
  =
  \begin{bmatrix} c_{a} \\ c_{b} \end{bmatrix}.
  \label{eq:lambda-multi}
\end{equation}
Los bloques de una misma variable son las sumas ya conocidas,
\[
  \Omega_{aa}(h) = \sum_{k} C^{aa}_{kk}(h),
  \qquad
  \Omega_{bb}(h) = \sum_{l} C^{bb}_{ll}(h),
\]
pero el bloque cruzado necesita los productos internos entre las funciones
propias de las dos variables, que ya no son ortogonales entre sí:
\begin{equation}
  \Omega_{ab}(h) = \sum_{k}\sum_{l} G_{kl}\; C^{ab}_{kl}(h),
  \qquad
  G_{kl} = \langle \xi^{a}_{k}, \xi^{b}_{l} \rangle .
  \label{eq:cruzado}
\end{equation}
Los lados derechos son $c_a = \Omega_{aa}(d_{a,0})$ y $c_b =
\Omega_{ba}(d_{b,0})$, y la varianza de predicción sería
\[
  \Omega_{aa}(0) - \begin{bmatrix} c_a \\ c_b \end{bmatrix}^{T}
  \begin{bmatrix}
    \Omega_{aa} & \Omega_{ab} \\
    \Omega_{ba} & \Omega_{bb}
  \end{bmatrix}^{-1}
  \begin{bmatrix} c_a \\ c_b \end{bmatrix}.
\]

\subsection{Cómo encaja en el \emph{solver}}

El sistema \eqref{eq:lambda-multi} tiene exactamente la forma que el
\emph{solver} ya resuelve. Lo que cambia es qué son los campos, las
covarianzas y los valores:

\begin{center}
\begin{tabular}{lll}
\hline
 & Método \code{scores} & Método \code{lambda} \\
\hline
Campo & un \emph{score} & una variable \\
Covarianza & $C_{ij}(h)$ del modelo & $\Omega_{ab}(h)$ agregada, ecuación \eqref{eq:cruzado} \\
Valores & un número por sitio & una curva por sitio (sus coeficientes) \\
Salida & \emph{scores} predichos & coeficientes de la curva predicha \\
\hline
\end{tabular}
\end{center}

El único cambio necesario en \code{simple\_kriging()} es aceptar que los
valores de cada campo sean una matriz (una fila por sitio, una columna por
coeficiente de la base) en lugar de un vector. Las fórmulas
\eqref{eq:triangular} valen sin cambio: $w = R^{-T} Z$ pasa a ser una matriz
y $V^{T} w$ da directamente los coeficientes de las curvas predichas.

Con ese cambio, el método univariado actual se podría reescribir sobre el
mismo \emph{solver} (un solo campo con covarianza $\sum_k C_{kk}$), lo que da
una prueba inmediata: debe reproducir los resultados de hoy.

\subsection{Lo que hace falta}

\begin{enumerate}
  \item \textbf{Un dominio y una base comunes.} Para sumar curvas de $b$ a la
        predicción de $a$, ambas deben estar definidas sobre el mismo
        intervalo de tiempo y expresadas en la misma base. Hoy no se exige:
        en los datos de México, PM10 tiene 4344 instantes y NO2 tiene 4292.
        Habría que exigirlo para este método o proyectar las curvas de una
        variable en la base de la otra.
  \item \textbf{La matriz $G$.} Con una base común de matriz de productos
        internos $J$, y con $\Xi^{a}$, $\Xi^{b}$ las matrices de coeficientes
        de las funciones propias, $G = (\Xi^{a})^{T} J\, \Xi^{b}$.
  \item \textbf{Las covarianzas agregadas.} Una función por pareja de
        variables, construida a partir de los variogramas compilados y de
        $G$. Para las parejas de \emph{scores} sin variograma cruzado el
        término se omite.
  \item \textbf{El \emph{solver} con valores matriciales}, como se describió.
  \item \textbf{Dos decisiones de criterio.}
        \begin{itemize}
          \item Centrar las curvas. Con varias variables no parece
                defendible sumar curvas sin centrar de una variable distinta;
                esto obligaría a decidir también qué hacer con el caso
                univariado, que hoy sigue a CRAN.
          \item Usar la meseta del modelo o el valor propio en la varianza.
        \end{itemize}
\end{enumerate}

\subsection{Cómo se verificaría}

Como no hay referencia externa, la verificación tendría que apoyarse en
propiedades:
\begin{itemize}
  \item Sin variogramas cruzados entre variables, o con $G = 0$, debe dar
        exactamente el \code{lambda} univariado de cada variable.
  \item Debe coincidir con la solución directa de \eqref{eq:lambda-multi}
        mediante \code{solve()}, como en las verificaciones de este
        documento.
  \item Con curvas centradas, predecir en un sitio observado debe devolver la
        curva observada.
  \item La validación cruzada debe poder compararlo con el cokriging por
        \emph{scores} sobre el mismo modelo.
\end{itemize}

% ======================================================================
\section{Limitaciones del \emph{solver}}

\begin{itemize}
  \item \textbf{Vecindario global.} Usa todos los sitios a la vez. El costo
        crece con el cubo del número de observaciones de cada grupo.
  \item \textbf{Sin memoria entre llamadas.} Cada llamada a \code{predict()}
        arma y factoriza $\Sigma$ de nuevo. En la validación cruzada eso
        significa una factorización por sitio y por modelo.
  \item \textbf{Validez del modelo.} Solo se detecta si la factorización de
        Cholesky falla.
  \item \textbf{Solo kriging simple.} No hay restricción de insesgamiento, de
        modo que los modelos sin meseta no se pueden usar.
  \item \textbf{Distancias euclidianas} sobre las coordenadas tal como vienen.
\end{itemize}

\end{document}
